\chapter{Propuesta}
\label{cap:propuesta}
Este capítulo detalla la metodología propuesta: un \textit{pipeline} liviano para \textit{Action Quality Assessment} basado en arquitecturas modernas de bajo costo computacional. La propuesta se centra en dos componentes esenciales: (i) la elección de \textit{Students} livianos modernos (\ac{TSM}-MobileNetV2 y MobileNetV3) inicializados con pesos preentrenados en ImageNet, y (ii) una receta de entrenamiento moderna (AdamW, programación cosenoidal, precisión mixta, acumulación de gradientes) coherente con el preprocesamiento del \textit{Teacher} \ac{I3D}. Sobre esta línea base se evalúa, como análisis marginal en el Capítulo~\ref{chap:resultados}, un esquema de destilación de conocimiento (\ac{KD}) con el objetivo de cuantificar su aporte real cuando la línea base ya es fuerte.

% ---
\section{Planteamiento de la solución}

La hipótesis del trabajo es que los \textit{Students} livianos preentrenados en ImageNet, al combinar un \textit{pipeline} moderno de entrenamiento con un preprocesamiento coherente con el \textit{Teacher} (mismas estadísticas de Kinetics-400, misma longitud de clip $T=64$ a 25~FPS, mismas transformaciones espaciales), pueden cubrir la mayor parte de la brecha de rendimiento respecto al \textit{Teacher} \ac{I3D} sin necesidad de mecanismos adicionales. La validación empírica de esta hipótesis es la contribución central de la tesis.

El \textit{pipeline} propuesto se compone de los siguientes elementos, descritos en detalle en las secciones siguientes:

\begin{enumerate}
    \item Arquitecturas \textit{Student} livianas (Sección~\ref{sec:modelos}).
    \item Pérdida de regresión $\mathcal{L}_{reg}$ sobre la etiqueta de calidad.
    \item Receta de entrenamiento moderna: AdamW, cosine annealing, AMP, gradient accumulation (Sección~\ref{sec:evaluacion}).
    \item \textit{Pipeline} de datos uniforme para los tres dominios evaluados (Sección~\ref{sec:datasets}).
    \item Protocolo de evaluación con métricas de precisión y eficiencia computacional (Sección~\ref{sec:evaluacion}).
\end{enumerate}

\subsection{Función de pérdida principal}

El \textit{Student} se entrena con la pérdida de regresión sobre la puntuación normalizada $y \in [0, 1]$:

\[
\mathcal{L}_{reg}(\hat{y}, y) = \frac{1}{N} \sum_{i=1}^{N} (\hat{y}_i - y_i)^2,
\]

donde $\hat{y}$ es la predicción del \textit{Student} y $y$ la etiqueta real. Esta pérdida es suficiente para sostener la contribución central del trabajo. Los Capítulos~\ref{chap:resultados} y la Sección~\ref{sec:kd_marginal} evalúan, como análisis complementario, el aporte de pérdidas adicionales de destilación.

\subsection{Análisis marginal de la destilación}
\label{sec:kd_marginal}

Para responder \textbf{OE6} sobre la necesidad real del \ac{KD} cuando se dispone de un \textit{pipeline} liviano fuerte, se considera además un esquema de destilación con dos pérdidas auxiliares aplicadas durante el entrenamiento del \textit{Student}:

\begin{itemize}
    \item \textbf{Destilación de atención espacio-temporal ($\mathcal{L}_{att}$)}: se extraen mapas de atención derivados de las activaciones convolucionales 3D del \textit{Teacher} (\ac{I3D}), de dimensiones $T \times H \times W$, y se alinean con las activaciones internas del \textit{Student} mediante una pérdida de divergencia (\ac{MSE} o \ac{KL}).
    \item \textbf{Alineación de representaciones intermedias ($\mathcal{L}_{temp}$)}: se proyectan a una dimensión común (usando un \textit{Conv3D} $1\times1\times1$) las activaciones de capas intermedias del \textit{Teacher} y del \textit{Student}, y se aplica \ac{MSE} cuadro a cuadro para forzar la sincronización temporal.
\end{itemize}

La pérdida total combinada queda como:

\[
\mathcal{L}_{total} = \alpha \cdot \mathcal{L}_{reg} + \beta \cdot \mathcal{L}_{att} + \gamma \cdot \mathcal{L}_{temp},
\]

donde la línea base del \textit{pipeline} liviano corresponde a $\beta = \gamma = 0$ (sólo $\mathcal{L}_{reg}$). Activar $\beta, \gamma > 0$ es la variante \textit{+~KD} que se evalúa como análisis marginal en el Capítulo~\ref{chap:resultados}. Durante el entrenamiento con \ac{KD}, el \textit{Teacher} permanece congelado y se utiliza únicamente para extraer activaciones de referencia.

Este diseño permite responder de forma directa la pregunta: \textit{¿qué aporta la destilación cuando la línea base liviana ya es fuerte?} Los resultados se reportan en la Sección~\ref{sec:kd_marginal_results}.


% ---
\section{Modelos seleccionados}
\label{sec:modelos}

El \textit{pipeline} liviano se construye sobre dos arquitecturas \textit{Student} contrastantes en cuanto a la presencia de modelado temporal explícito. Como referencias de capacidad alta se incluyen además dos \textit{Teachers} 3D: \ac{I3D} (referencia histórica del campo) y SlowFast-R50 (referencia 3D moderna). Esta combinación permite cuantificar la brecha del \textit{pipeline} liviano frente a dos paradigmas 3D distintos y evaluar el aporte marginal de la destilación sobre una línea base fuerte.

\subsection{Modelos Teacher (referencia)}

Los \textit{Teachers} son modelos 3D profundos preentrenados en \textit{Kinetics-400}. En esta tesis cumplen únicamente el rol de referencia de alta capacidad para cuantificar la brecha del \textit{pipeline} liviano (\textbf{OE3}); no son la propuesta. Se utilizan dos:

\begin{itemize}
    \item \textbf{I3D-R50} (\textit{Inflated 3D ConvNet}, \cite{Carreira2017I3D}): basado en la expansión de filtros 2D preentrenados a 3D. Es la referencia histórica del campo \ac{AQA} y permite comparabilidad con la literatura previa.
    \item \textbf{SlowFast-R50} (\cite{feichtenhofer2019slowfast}): arquitectura 3D moderna con dos \textit{pathways} (uno lento de alta resolución temporal espacial, otro rápido de baja resolución pero alta tasa de muestreo). Permite cuantificar dónde se sitúa el \textit{pipeline} liviano respecto al paradigma 3D actual, no sólo al histórico. Sólo se entrena en AQA-7 (dataset principal del trabajo).
\end{itemize}

Ambos \textit{Teachers} se entrenan con la misma pérdida de regresión $\mathcal{L}_{reg}$ que los \textit{Students}, sin destilación, para mantener la comparación limpia.

\subsection{Modelos Student (pipeline liviano propuesto)}

Los \textit{Students} son la contribución central de la tesis. Se seleccionan dos arquitecturas livianas con distinto grado de modelado temporal explícito, lo que permite analizar (en la ablación E7 del Capítulo~\ref{chap:resultados}) qué tanto aporta el componente temporal en \ac{AQA}.

\begin{itemize}
    \item \textbf{TSM-MobileNetV2} \cite{Lin2019TSM,sandler2018mobilenetv2}:
        \begin{itemize}
            \item MobileNetV2 adaptado a video mediante \textit{Temporal Shift Modules} (\ac{TSM}), que desplazan parte de los canales de cada \textit{frame} en el tiempo. Permite modelar relaciones temporales sin convoluciones 3D.
            \item Reutiliza pesos preentrenados de ImageNet con coste computacional adicional mínimo.
            \item Es la arquitectura principal del \textit{pipeline} propuesto.
        \end{itemize}
    \item \textbf{MobileNetV3-Large} \cite{Howard2019MobileNetV3}:
        \begin{itemize}
            \item Arquitectura optimizada para eficiencia, con activaciones \textit{h-swish} y bloques \textit{Squeeze-and-Excitation}.
            \item Aplicada a video mediante procesamiento por \textit{frame} con agregación temporal por \textit{pooling} global, sin mecanismos explícitos de modelado temporal.
            \item Su inclusión permite contrastar el aporte del componente temporal explícito.
        \end{itemize}
\end{itemize}

Cada \textit{Student} se entrena con el \textit{pipeline} liviano propuesto (sólo $\mathcal{L}_{reg}$, sin participación del \textit{Teacher}). Adicionalmente, como análisis marginal del Capítulo~\ref{chap:resultados}, se reporta la variante \textit{+~KD} que incorpora $\mathcal{L}_{att}$ y $\mathcal{L}_{temp}$ desde el \textit{Teacher} \ac{I3D}, con el objetivo de cuantificar si la destilación aporta valor sobre la línea base liviana o si resulta innecesaria.


% ---
\section{Datasets y preprocesamiento}
\label{sec:datasets}

\subsection{Conjuntos de datos}

Se emplearon tres conjuntos de datos ampliamente utilizados en la tarea de \textit{Action Quality Assessment} (AQA). Cada uno posee características distintas en cuanto a dominio, número de muestras, estructura y tipo de anotaciones. A continuación se detallan:

Los tres conjuntos son distribuidos como \textit{frames} pre-extraídos en directorios (no como archivos de video), lo cual simplifica el pipeline de preprocesamiento al eliminar la dependencia de decodificadores de video y permitir un acceso directo por frame.

\begin{itemize}
    \item \textbf{AQA-7}:
        \begin{itemize}
            \item \textbf{Cantidad de muestras}: 1,189 clips distribuidos en 7 disciplinas deportivas distintas: clavados, gimnasia artística, salto con pértiga, esquí estilo libre, patinaje artístico, patinaje en línea y snowboard.
            \item \textbf{Duración promedio por clip}: 5 a 8 segundos.
            \item \textbf{Resolución promedio}: 720p.
            \item \textbf{Peso total del dataset}: Aproximadamente 3.6 GB.
            \item \textbf{Anotaciones}: Etiquetas de puntuación continua específicas para cada disciplina, normalmente en rangos de 0 a 100. Se evalúan mediante métricas SRCC, PLCC y MAE.
            \item \textbf{Formato}: directorios con frames en .jpg ordenados por índice temporal.
            \item \textbf{Rol en la tesis}: dataset principal multi-disciplina; soporte de la Tabla de resultados principal y punto central de la generalización cruzada.
        \end{itemize}

    \item \textbf{MTL-AQA}:
        \begin{itemize}
            \item \textbf{Cantidad de muestras}: 1{,}412 clips de clavados desde trampolín y plataforma.
            \item \textbf{Duración promedio por clip}: 3.5 segundos ($\sim$109 frames a 30~FPS).
            \item \textbf{Resolución promedio}: 720p.
            \item \textbf{Peso total del dataset}: Aproximadamente 2.8 GB.
            \item \textbf{Anotaciones}: el archivo oficial \texttt{final\_annotations\_dict.pkl} provee, por clip, una puntuación final (\texttt{final\_score}) además de subcomponentes (dificultad, posición, tipo de rotación, etc.). En esta tesis se utiliza únicamente \texttt{final\_score} como etiqueta de regresión; los subcomponentes no se emplean.
            \item \textbf{Formato}: directorios \texttt{<dive\_id>\_<clip\_id>/} con frames en \texttt{.jpg}.
            \item \textbf{Rol en la tesis}: corpus especializado en clavados utilizado para entrenar Teacher y Students; también funciona como dominio \textit{source} en el experimento de generalización cruzada (entrenamiento en MTL-AQA, evaluación en AQA-7).
        \end{itemize}

    \item \textbf{JIGSAWS}:
        \begin{itemize}
            \item \textbf{Cantidad de muestras}: 103 \textit{trials} oficiales de tres tareas quirúrgicas básicas (\textit{Suturing}, \textit{Knot Tying}, \textit{Needle Passing}) realizadas en el sistema robótico \textit{da Vinci}, distribuidos en dos vistas (\texttt{capture1} y \texttt{capture2}) para un total de 206 clips.
            \item \textbf{Duración promedio por trial}: 1 a 2 minutos. Para alimentar al modelo se muestrean 64 frames uniformemente distribuidos a lo largo del trial; no se realiza segmentación por gesto.
            \item \textbf{Resolución promedio}: 640$\times$480.
            \item \textbf{Peso total del dataset}: Aproximadamente 1.5 GB en formato frames.
            \item \textbf{Anotaciones}: puntuaciones OSATS (\textit{Objective Structured Assessment of Technical Skills}) compuestas por seis ítems en el rango $[1,5]$ que suman GRS $\in [6,30]$, además de etiquetas de \textit{expertise} (\textit{novice}, \textit{intermediate}, \textit{expert}). Se utiliza GRS como etiqueta de regresión.
            \item \textbf{Formato}: directorios \texttt{<Tarea>\_<Sujeto\_Trial>\_capture<N>/} con frames en \texttt{.jpg}.
            \item \textbf{Rol en la tesis}: dataset fuera del dominio deportivo que amplía la evaluación del \textit{framework} al campo de la cirugía robótica, permitiendo verificar la robustez cross-dominio del modelo.
        \end{itemize}
\end{itemize}


\subsection{Preprocesamiento}
El pipeline de preprocesamiento se divide en cuatro etapas:

\subsubsection{Extracción y muestreo de frames}
\begin{itemize}
    \item \textbf{Herramienta}: \textit{OpenCV} (cv2.VideoCapture).
    \item \textbf{Tasa de muestreo}: 25 FPS (balance entre fluidez temporal y costo computacional).
    \item \textbf{Filtrado}: Eliminación de frames corruptos o vacíos mediante umbralización de histograma.
\end{itemize}

\subsubsection{Segmentación temporal}
\begin{itemize}
    \item \textbf{Duración de clips}: 64 frames por video o trial en los tres datasets. A 25 FPS esto corresponde a 2.56 segundos; cuando el material original viene a 30~FPS efectivos (MTL-AQA, JIGSAWS) los frames se obtienen mediante muestreo uniforme a lo largo de la secuencia para conservar el tamaño temporal $T = 64$.
    \item \textbf{Solapamiento}: cada video produce un único clip de 64 frames; no se generan clips solapados.
    \item \textbf{Alineación con anotaciones}: cada clip hereda directamente la puntuación final del video o trial original (\texttt{final\_score} en MTL-AQA, GRS en JIGSAWS, score por disciplina en AQA-7).
\end{itemize}

\subsubsection{Transformación espacial}
\begin{itemize}
    \item \textbf{Redimensionamiento}: Interpolación bicúbica a $224 \times 224$ píxeles (compatible con las arquitecturas base).
    \item \textbf{Normalización}: Ajuste de valores de píxeles usando estadísticos de \textit{Kinetics-400}:
        \[
        \text{Pixel}_{\text{norm}} = \frac{\text{Pixel}_{\text{raw}} - [0.43216, 0.394666, 0.37645]}{[0.22803, 0.22145, 0.216989]}
        \]
        donde los valores corresponden a (media, std) por canal RGB.
    \item \textbf{Aumento de datos}: Durante entrenamiento, se aplican:
        \begin{itemize}
            \item Volteo horizontal aleatorio (probabilidad 0.5).
            \item Recorte temporal aleatorio (variación de $\pm10\%$ en la duración del clip).
        \end{itemize}
\end{itemize}

\subsubsection{Formateo final}
\begin{itemize}
    \item \textbf{Estructura}: Tensores de dimensión $[B, C, T, H, W]$, donde:
        \begin{itemize}
            \item $B$: Tamaño del \textit{batch} efectivo. Por restricciones de VRAM en la GPU utilizada (RTX 3060 Mobile, 6~GB), el batch físico se fija en 2 para Students y 1 para entrenamientos con KD; el batch efectivo de 16 se obtiene con \textit{gradient accumulation} de 8 o 16 pasos respectivamente.
            \item $C$: Canales (3 para RGB).
            \item $T$: Frames temporales ($T = 64$ uniforme en todos los experimentos).
            \item $H, W$: Altura y anchura (224).
        \end{itemize}
    \item \textbf{Compatibilidad}: Los tensores se ajustan a los requerimientos de entrada de I3D (Teacher) y TSM-MobileNetV2 / MobileNetV3 (Students).
\end{itemize}

\subsection{Consideraciones técnicas}
\begin{itemize}
    \item \textbf{Balanceo}: para datasets con distribución sesgada (MTL-AQA con concentración en puntuaciones medias y JIGSAWS con clases desbalanceadas por tarea y nivel de \textit{expertise}) se aplica muestreo estratificado durante la construcción de los splits de entrenamiento y validación.
    \item \textbf{División}: para AQA-7 se respeta el split oficial \textit{Split\_4} (1106 clips, train/test predefinido) y el conjunto de validación se obtiene tomando el 15\% del train oficial mediante muestreo estratificado por categoría. Para MTL-AQA y JIGSAWS, donde no existe un split oficial canónico, se aplica una división estratificada 70\%/15\%/15\% por (categoría $\times$ bin de score) usando \texttt{StratifiedShuffleSplit}.
    \item \textbf{Reproducibilidad}: semillas fijas para inicialización, \textit{shuffling} y muestreo de frames. Todos los resultados reportados corresponden a la semilla 42; la extensión a múltiples semillas (con reporte de media y desviación estándar) se identifica como trabajo futuro.
\end{itemize}

% ---
\section{Metodología de Evaluación}
\label{sec:evaluacion}

\subsection{Métricas de Precisión}

Para cuantificar el rendimiento de los modelos en la tarea de \textit{Action Quality Assessment} (AQA), se emplearán las siguientes métricas estadísticas:

\begin{itemize}
    \item \textbf{Coeficiente de Correlación de Spearman (SRCC)}: Evalúa la correlación monotónica entre los rangos de las predicciones del modelo y los rangos de las puntuaciones reales. Es útil cuando no se asume una relación lineal exacta entre las variables.
    \begin{itemize}
        \item \textit{Unidad}: Adimensional, valor en el rango \([-1, 1]\).
        \item \textit{Interpretación}: Un valor cercano a \(1\) indica una alta concordancia en el orden relativo de las predicciones y etiquetas. \textbf{Un valor más alto es mejor}.
    \end{itemize}
    \[
    \rho = 1 - \frac{6 \sum d_i^2}{n(n^2 - 1)}
    \]
    donde \(d_i\) es la diferencia entre los rangos de las predicciones y las etiquetas para cada muestra, y \(n\) es el número total de muestras.

    \item \textbf{Coeficiente de Correlación Lineal de Pearson (PLCC)}: Mide la relación lineal entre los valores predichos y las etiquetas reales.
    \begin{itemize}
        \item \textit{Unidad}: Adimensional, valor en el rango \([-1, 1]\).
        \item \textit{Interpretación}: Un valor cercano a \(1\) indica una fuerte relación lineal positiva. \textbf{Un valor más alto es mejor}.
    \end{itemize}
    \[
    r = \frac{\sum (x_i - \bar{x})(y_i - \bar{y})}{\sqrt{\sum (x_i - \bar{x})^2 \sum (y_i - \bar{y})^2}}
    \]

    \item \textbf{Error Absoluto Medio (MAE)}: Mide el error promedio absoluto entre las predicciones del modelo y las puntuaciones reales, en la misma escala original de las etiquetas.
    \begin{itemize}
        \item \textit{Unidad}: La misma que las puntuaciones (por ejemplo, puntos sobre 100).
        \item \textit{Interpretación}: Cuanto menor sea el MAE, mayor será la precisión del modelo. \textbf{Un valor más bajo es mejor}.
    \end{itemize}
    \[
    \text{MAE} = \frac{1}{n} \sum_{i=1}^n |y_i - \hat{y}_i|
    \]
\end{itemize}


\subsection{Métricas de Eficiencia}

Para evaluar la viabilidad computacional de los modelos propuestos, se utilizarán las siguientes métricas cuantitativas:

\begin{itemize}
    \item \textbf{FLOPs (Floating Point Operations)}: Representa el número total de operaciones de punto flotante necesarias para procesar un clip de entrada. Se calcula utilizando la librería \texttt{fvcore}.
    \begin{itemize}
        \item \textit{Unidad}: GigaFLOPs (GFLOPs).
        \item \textit{Interpretación}: Menores valores indican menor complejidad computacional. \textbf{Un valor más bajo es mejor}, ya que implica menor consumo de recursos y mayor eficiencia.
    \end{itemize}

    \item \textbf{Latencia de Inferencia}: Tiempo promedio que el modelo requiere para procesar un clip de entrada, medido en la GPU utilizada en este trabajo (NVIDIA RTX 3060 Mobile, 6~GB de VRAM). Este valor excluye el tiempo de carga de datos y refleja únicamente el tiempo de ejecución del modelo.
    \begin{itemize}
        \item \textit{Unidad}: Milisegundos (ms).
        \item \textit{Interpretación}: Una latencia baja es deseable para aplicaciones en tiempo real. \textbf{Un valor más bajo es mejor}, ya que indica mayor rapidez de respuesta.
    \end{itemize}
\end{itemize}

Estas métricas permiten evaluar no solo la precisión, sino también la factibilidad de implementar los modelos en entornos reales, particularmente en dispositivos de bajo consumo o aplicaciones que requieren respuestas inmediatas.


\subsection{Protocolo Experimental}
\begin{enumerate}
    \item \textbf{División de Datos}: para AQA-7 se respeta el \textit{Split\_4} oficial; para MTL-AQA y JIGSAWS se aplica un split estratificado 70/15/15 por (categoría $\times$ bin de score). Las proporciones por disciplina o tarea se mantienen entre conjuntos.

    \item \textbf{Comparativas en el dataset principal (AQA-7)}: se reportan los siguientes modelos:
    \begin{itemize}
        \item \textbf{Teachers de referencia}: I3D-R50 y SlowFast-R50, ambos entrenados con $\mathcal{L}_{reg}$.
        \item \textbf{\textit{Pipeline} liviano (propuesto)}: TSM-MobileNetV2 y MobileNetV3 con $\mathcal{L}_{reg}$, evaluados sobre tres semillas (42, 0, 7) para reportar media y desviación estándar.
        \item \textbf{Ablaciones del \textit{pipeline}}: TSM-MobileNetV2 sin preentrenamiento ImageNet (E6) y MobileNetV2 plano sin \ac{TSM} (E7), ambos con semilla 42.
        \item \textbf{Análisis marginal de la destilación (\textit{+~KD})}: TSM-MobileNetV2 y MobileNetV3 entrenados con $\mathcal{L}_{total}$ completo ($\beta, \gamma > 0$), con semilla 42.
    \end{itemize}

    \item \textbf{Comparativas en MTL-AQA y JIGSAWS}: se reportan I3D \textit{Teacher}, ambos \textit{Students} con $\mathcal{L}_{reg}$ y la variante \textit{+~KD} para análisis marginal, con semilla 42.

    \item \textbf{Repetibilidad}: todos los resultados se obtienen con \textit{seeds} explícitas para inicialización, \textit{shuffling}, muestreo y \textit{augmentation}. En el dataset principal (AQA-7) los \textit{Students} se ejecutan con tres semillas; en los demás datasets se utiliza la semilla 42 para limitar el costo computacional acumulado.
\end{enumerate}

\subsection{Análisis Adicionales}
\begin{itemize}
    \item \textbf{Mapas de Atención}: visualización de las regiones espaciotemporales consideradas críticas por los modelos, generadas con Grad-CAM sobre el bloque intermedio de cada arquitectura. Se reportan ejemplos representativos en los tres datasets (Sección~\ref{sec:gradcam}).

    \item \textbf{Evaluación Cruzada}: entrenamiento en un dominio \textit{source} y evaluación \textit{zero-shot} sobre el conjunto de test de un dominio \textit{target} distinto. Se consideran dos transferencias: MTL-AQA $\to$ AQA-7 (dominios deportivos relacionados) y AQA-7 $\to$ JIGSAWS (deporte $\to$ cirugía robótica).

    \item \textbf{Benchmark de Hardware}: medición de FLOPs (vía \texttt{fvcore}) y latencia mediana en RTX 3060 Mobile como aproximación a hardware con recursos limitados. La validación en plataformas embebidas específicas (Jetson Nano, móviles, etc.) se identifica como trabajo futuro.
\end{itemize}